% Subsection 2.3.1: Zero-shot CoT

零样本链式思维（Zero-shot CoT）不依赖任何示例，仅通过简单指令触发模型自发的分步推理，其核心是利用大语言模型自身的泛化能力完成步骤化推导\cite{kojima2022large,wei2022chain,li2024cotkd}。典型形式是在输入提示中加入类似“请一步步思考”的引导语句，使模型将单阶段生成过程拆分为两阶段：
\[
\text{问题} \rightarrow \text{推理步骤} \rightarrow \text{代码结果}
\]
Zero-shot CoT 具有实现简单、无额外数据成本、不占用上下文空间的优势，能够在不改变模型参数的情况下提升基础逻辑表达能力，常作为推理增强的基线方案。在算法代码生成中，它可强制模型先明确算法思路、边界处理与步骤规划，再输出代码，从而降低直接生成带来的逻辑偏差。